\chapter{Simkin's Theorem and the Queenon}
\label{ch:simkin}

Chapter~\ref{ch:gauss} stated Simkin\index{Simkin's theorem}'s theorem that $Q(N) = ((1 \pm o(1)) N e^{-\alpha})^N$ with $\alpha \approx 1.942$. This chapter describes how the theorem is proved and what the constant means, because the structure of the proof is itself a vertical argument: a lower bound from a sequential construction that is a ladder, and an upper bound from an entropy count that is a width.

\section{The upper bound: entropy}

Let $X$ be a uniformly random solution of the $N$ queens problem, encoded as the permutation $(x_1, \dots, x_N)$. Its entropy is $H(X) = \log Q(N)$, and the chain rule gives
\[
\log Q(N) = \sum_{i=1}^{N} H(x_i \mid x_1, \dots, x_{i-1}) .
\]
Each conditional entropy is at most the logarithm of the number of columns available to the queen in row $i$ given the earlier rows, and a column is available only if it is not already used and not attacked diagonally by any earlier queen. The trivial count, at most $N - (i-1)$ unused columns in row $i$, gives only $Q(N) \le N!$, that is $Q(N) \le (N e^{-1})^N$ up to a polynomial factor. This already gives the right form $(cN)^N$ but not the constant, since the truth is $c = e^{-\alpha}$ with $\alpha \approx 1.942$ rather than $1$. The diagonal constraints cannot be handled by subtracting attacked columns row by row, because the attacked columns overlap in a way that depends on the earlier placements.

Simkin's upper bound refines the count by conditioning on the empirical distribution of the solution on a coarse grid. Divide the board into $k \times k$ blocks and let $\nu$ be the matrix of block occupation numbers, which sums to $N$ and has row and column sums $N/k$. Given $\nu$, the entropy of $X$ is at most the entropy of a uniformly random placement with those block counts that respects the diagonal constraints within and between blocks, and this can be bounded by a sum over blocks of a local entropy term plus a correction for the diagonals. Letting $k \to \infty$ slowly with $N$, the bound becomes a functional of the limiting density $f(x, y)$ of queens on the unit square, and the best upper bound is the supremum of that functional over admissible $f$. The admissible limits are the queenon\index{queenon}s: permutons (probability measures on $[0,1]^2$ with uniform marginals in each coordinate, not necessarily with densities) that in addition satisfy the diagonal constraint that the mass on every diagonal and antidiagonal strip is at most the width of the strip, measured in the coordinate $x + y$ or $x - y$. For a queenon with density $f$ the functional is schematically
\[
\Psi(f) = -\int_{[0,1]^2} f \log f + \text{(diagonal terms)} ,
\]
where the first term is the entropy and the diagonal terms account for the requirement that each diagonal carry at most one queen. The best upper bound is the supremum of this functional over queenons, and the constant $\alpha$ is determined by that supremum. We do not reproduce Simkin's exact formula here (see \cite{Simkin}), and his numerical optimisation gives $\alpha = 1.942 \pm 0.003$.

\section{The lower bound: the random greedy ladder}

The matching lower bound is constructive. Place the queens one at a time in a random order of rows, choosing at each step a uniformly random column among those not attacked by the queens already placed, and analyse the number of available columns as the process runs. This is a random greedy algorithm, and Simkin shows that with the right choice of the queenon as a target distribution for the placements, and a careful analysis of the differential equations that govern the expected number of available positions, the algorithm succeeds in placing all but a vanishing fraction of the queens with probability bounded below, and the number of distinct outcomes is at least $((1 - o(1)) N e^{-\alpha})^N$. The argument uses the differential equations method of Wormald and an absorption step to complete the last few queens.

In our terms the random greedy algorithm is a ladder. Each row is a rung, the available columns at row $i$ are the level set $\{ x_i : S(x_1, \dots, x_i) = 0 \}$ given the earlier placements, and the number of outcomes is the product over rungs of the number of available choices. The lower bound is the product of the widths along the ladder, and the queenon is the choice of weight that makes the product as large as possible. This is the self reducibility route of Jerrum, Valiant\index{Valiant, Leslie} and Vazirani\index{Jerrum, Valiant and Vazirani}, row by row, and the differential equations are a deterministic approximation to the vertical walk in the limit $N \to \infty$. Simkin's proof is a derandomized vertical walk for the limiting problem.

\section{The queenon}

The optimising density $f^*$ is not uniform. Our own enumeration of the solutions for $N = 13$ (\texttt{code/queenon\_output.txt}, density normalised so that the uniform value is $1$) shows the density largest near the midpoints of the edges (about $1.42$) and below uniform both at the centre of the board (about $0.79$) and in the corners (about $0.49$). We do not offer an explanation of this pattern. A uniform solution is, in the limit, a sample from $f^*$ in the sense that its empirical distribution on any coarse grid converges to the block integrals of $f^*$.

This is a large deviations statement. The uniform distribution on permutations has, by Sanov's theorem and its permutation analogue, a limiting empirical density that is uniform on the square; conditioning on the rare event $S = 0$ tilts the limit to $f^*$. The tilt accounts for only part of the rate $\alpha - 1$: the remainder is the diagonal (no collision) term, which is paid even at the uniform density, as the independence estimate below suggests. The queenon is to the queens what the Wigner semicircle is to random matrices or the Vershik Kerov shape is to random Young diagrams: the limit shape under a hard constraint.

The queenon is not the optimal importance density for the vertical walk, which is the conditional law $p(\cdot \mid S = 0)$ of Chapter~\ref{ch:rareevent}; its one point marginals are only precursors of the queenon. A sampler that could draw permutations whose empirical density approximates $f^*$ would remove the part of the cost that is due to the tilt, but not the diagonal part, so the ground step would still have exponentially small probability. The cross entropy\index{cross entropy method} method of Chapter~\ref{ch:rareevent} with a product proposal having these marginals reduces the relative second moment of crude Monte Carlo by a factor of at most $e^{d}$, at most about $2$ for $7 \le N \le 12$ ($e^{0.68}$ at $N = 8$, while $N = 6$ gives $e^{2.43}$) (Chapter~\ref{ch:ce}), and queenon tilted swaps do not speed up the bottom rung chain (Chapter~\ref{ch:open}).

\section{The constant}

What is $\alpha$? It is defined by a variational problem and computed numerically to three decimals. No closed form is known and none is expected; the variational problem is infinite dimensional and its optimiser is not expressible in elementary terms. Simkin's bounds are $1.939 \le \alpha \le 1.945$. For comparison, if the pairs were independent with the mean of Chapter~\ref{ch:phase} the constant would be $1 + 2/3 = 1.667$; the gap of $0.275$ is the correlation correction.

For the walk the constant sets the length of the ladder: $\log(1/\rho_N) = (\alpha - 1) N + o(N) \approx 0.942 N$ on permutations, so with ratio $\tfrac12$ per rung one needs $1.36 N$ rungs. The experiments of Chapter~\ref{ch:experiments} used between five and nine rungs for $8 \le N \le 16$. Since the ratios along a ladder multiply to $\rho_N$, this forces a small average ratio: for $N = 16$, $\log(16!/Q(16)) \approx 14.2$ over nine rungs is an average ratio of about $0.21$, made up of ratios of $0.3$ to $0.6$ at the upper rungs and much smaller ones at the bottom (the final ratio is about $0.01$ at $N = 16$).

\section{Finite size corrections}

The $o(1)$ in Simkin's theorem is not small for the $N$ at which one can compute. From Table~\ref{tab:counts}, $\log_{10}(Q(N)/N!)$ decreases by $0.41$ per unit of $N$ above $N = 15$, matching $(\alpha - 1)/\ln 10 = 0.409$, but the intercept is off: the exact values are larger than $e^{-0.942 N}$ by a nearly constant factor, about $2.4$ for $15 \le N \le 27$. A natural guess is that $Q(N) / N! \sim C N^{\beta} e^{-(\alpha - 1) N}$ for constants $C$ and $\beta$, by analogy with the asymptotics of Latin squares and other constrained permutation counts, and the data up to $N = 27$ are consistent with $\beta = 0$ (Chapter~\ref{ch:open}). Establishing the polynomial correction would require a second order large deviations analysis around the queenon, which has not been done, and we record it among the open problems.
